% GENERATED FILE. Edit canonical lessons, then run npm run generate:books.
% canonical-source-hash: fnv1a64:00848eb8feab5578
% canonical-lessons: RU-C452-bulochnaya, RU-C452-gastronom, RU-C452-prilavok, RU-C452-prodavshchitsa, RU-C452-vzveshivat

\chapter{Food shopping}
\label{ch:ru-food-shopping}
\hlchaptermodality{452}

\begin{chapteropening}
\textbf{By the end of this chapter:} \emph{I can shop for food: the bakery, the grocery, the counter, the shop assistant and weighing.}

Complete the last lesson of chapter 452: I can shop for food: the bakery, the grocery, the counter, the shop assistant and weighing.
\end{chapteropening}

\section[\texorpdfstring{\ru{булочная} (búlochnaya)}{булочная (búlochnaya)}]{\ru{булочная} (búlochnaya) — a bakery (shop)}

\label{lesson:RU-C452-bulochnaya}



\begin{quote}
Before the new one: say the Russian for an addressee, then the Russian for to issue.
\end{quote}

\subsection*{You'll want to know: \ru{булочная}}
\textbf{\ru{булочная}} (\emph{búlochnaya}) — \textquotedblleft{}a bakery (shop)\textquotedblright{}.

\begin{quote}
\ru{Утром} \ru{я} \ru{покупаю} \ru{хлеб} \ru{в} \ru{булочной}.
\end{quote}

\begin{quote}
\emph{Útrom ya pokupáyu khleb v búlochnoy.}
\end{quote}

\begin{quote}
In the morning I buy bread at the bakery.
\end{quote}

An adjective used as a noun. Traditionally \emph{\ru{чн}} is pronounced \emph{\ru{шн}} here.

\begin{cousinweb}
\emph{\ru{булочная}} is built on \emph{\ru{булочка}}, \emph{\ru{булка}} 'a roll' with an adjective ending: the shop for rolls. EtymDB lists \emph{\ru{булка}} without a deeper source.
\end{cousinweb}

\begin{grammarlens}[title={What you've built}]
One more word for telling, explaining and describing at B1.
\end{grammarlens}

\subsection*{Your turn}
\begin{itemize}
\raggedright
  \item \emph{Say it:} \emph{búlochnaya}
  \item \emph{Say it:} \emph{búlochnaya}, once more
  \item \emph{Say it:} a sentence of your own with \emph{búlochnaya}
  \item \emph{Recall:} say the Russian for an addressee, then the Russian for to issue, then say \emph{búlochnaya} again
\end{itemize}

\begin{tcolorbox}[breakable,colback=teal!4,colframe=teal!35!black,title={Before you move on}]
What does \ru{булочная} mean? (\textquotedblleft{}A bakery (shop)\textquotedblright{}.) Say it once more.
\end{tcolorbox}
\section[\texorpdfstring{\ru{гастроном} (gastronóm)}{гастроном (gastronóm)}]{\ru{гастроном} (gastronóm) — a grocery store, a food shop}

\label{lesson:RU-C452-gastronom}



\begin{quote}
Before the new one: say the Russian for to issue, then the Russian for a bakery.
\end{quote}

\subsection*{You'll want to know: \ru{гастроном}}
\textbf{\ru{гастроном}} (\emph{gastronóm}) — \textquotedblleft{}a grocery store, a food shop\textquotedblright{}.

\begin{quote}
\ru{В} \ru{гастрономе} \ru{на} \ru{углу} \ru{мы} \ru{купили} \ru{сыр} \ru{и} \ru{колбасу}.
\end{quote}

\begin{quote}
\emph{V gastronóme na uglú my kupíli syr i kolbasú.}
\end{quote}

\begin{quote}
We bought cheese and sausage at the grocery store on the corner.
\end{quote}

Masculine. Also an old word for a gourmet.

\begin{cousinweb}
\emph{\ru{гастроном}} is usually traced, probably through French \emph{gastronome}, to Greek \emph{gast\'{\={e}}r} 'belly' and \emph{nómos} 'law': first a gourmet, then a food shop. EtymDB does not list it, so the line is hedged.
\end{cousinweb}

\begin{grammarlens}[title={What you've built}]
One more word for telling, explaining and describing at B1.
\end{grammarlens}

\subsection*{Your turn}
\begin{itemize}
\raggedright
  \item \emph{Say it:} \emph{gastronóm}
  \item \emph{Say it:} \emph{gastronóm}, once more
  \item \emph{Say it:} a sentence of your own with \emph{gastronóm}
  \item \emph{Recall:} say the Russian for to issue, then the Russian for a bakery, then say \emph{gastronóm} again
\end{itemize}

\begin{tcolorbox}[breakable,colback=teal!4,colframe=teal!35!black,title={Before you move on}]
What does \ru{гастроном} mean? (\textquotedblleft{}A grocery store, a food shop\textquotedblright{}.) Say it once more.
\end{tcolorbox}
\section[\texorpdfstring{\ru{прилавок} (prilávok)}{прилавок (prilávok)}]{\ru{прилавок} (prilávok) — a counter (in a shop)}

\label{lesson:RU-C452-prilavok}



\begin{quote}
Before the new one: say the Russian for a bakery, then the Russian for a grocery store.
\end{quote}

\subsection*{You'll want to know: \ru{прилавок}}
\textbf{\ru{прилавок}} (\emph{prilávok}) — \textquotedblleft{}a counter (in a shop)\textquotedblright{}.

\begin{quote}
\ru{Продавец} \ru{стоит} \ru{за} \ru{прилавком}.
\end{quote}

\begin{quote}
\emph{Prodavéts stoít za prilávkom.}
\end{quote}

\begin{quote}
The shop assistant stands behind the counter.
\end{quote}

Masculine; the \ru{о} drops out: \emph{\ru{прилавка}}.

\begin{cousinweb}
\emph{\ru{прилавок}} is built from \emph{\ru{при}-} 'beside' and \emph{\ru{лавка}} 'a bench; a shop', which is inherited from Proto-Slavic \emph{lavŭka} 'a bench, a shop' (EtymDB).
\end{cousinweb}

\begin{grammarlens}[title={What you've built}]
One more word for telling, explaining and describing at B1.
\end{grammarlens}

\subsection*{Your turn}
\begin{itemize}
\raggedright
  \item \emph{Say it:} \emph{prilávok}
  \item \emph{Say it:} \emph{prilávok}, once more
  \item \emph{Say it:} a sentence of your own with \emph{prilávok}
  \item \emph{Recall:} say the Russian for a bakery, then the Russian for a grocery store, then say \emph{prilávok} again
\end{itemize}

\begin{tcolorbox}[breakable,colback=teal!4,colframe=teal!35!black,title={Before you move on}]
What does \ru{прилавок} mean? (\textquotedblleft{}A counter (in a shop)\textquotedblright{}.) Say it once more.
\end{tcolorbox}
\section[\texorpdfstring{\ru{продавщица} (prodavshchítsa)}{продавщица (prodavshchítsa)}]{\ru{продавщица} (prodavshchítsa) — a shop assistant (woman)}

\label{lesson:RU-C452-prodavshchitsa}



\begin{quote}
Before the new one: say the Russian for a grocery store, then the Russian for a counter.
\end{quote}

\subsection*{You'll want to know: \ru{продавщица}}
\textbf{\ru{продавщица}} (\emph{prodavshchítsa}) — \textquotedblleft{}a shop assistant (woman)\textquotedblright{}.

\begin{quote}
\ru{Продавщица} \ru{улыбнулась} \ru{и} \ru{дала} \ru{мне} \ru{сдачу}.
\end{quote}

\begin{quote}
\emph{Prodavshchítsa ulybnúlas i dalá mne sdáchu.}
\end{quote}

\begin{quote}
The shop assistant smiled and gave me my change.
\end{quote}

The feminine of \emph{\ru{продавец}}.

\begin{cousinweb}
\emph{\ru{продавщица}} is the feminine of \emph{\ru{продавец}} 'a seller', from \emph{\ru{продавать}} 'to sell', on \emph{\ru{давать}} 'to give', which is inherited from Proto-Slavic \emph{dāvati} (EtymDB).
\end{cousinweb}

\begin{grammarlens}[title={What you've built}]
One more word for telling, explaining and describing at B1.
\end{grammarlens}

\subsection*{Your turn}
\begin{itemize}
\raggedright
  \item \emph{Say it:} \emph{prodavshchítsa}
  \item \emph{Say it:} \emph{prodavshchítsa}, once more
  \item \emph{Say it:} a sentence of your own with \emph{prodavshchítsa}
  \item \emph{Recall:} say the Russian for a grocery store, then the Russian for a counter, then say \emph{prodavshchítsa} again
\end{itemize}

\begin{tcolorbox}[breakable,colback=teal!4,colframe=teal!35!black,title={Before you move on}]
What does \ru{продавщица} mean? (\textquotedblleft{}A shop assistant (woman)\textquotedblright{}.) Say it once more.
\end{tcolorbox}
\section[\texorpdfstring{\ru{взвешивать} (vzvéshivat)}{взвешивать (vzvéshivat)}]{\ru{взвешивать} (vzvéshivat) — to weigh}

\label{lesson:RU-C452-vzveshivat}



\begin{quote}
Before the new one: say the Russian for a counter, then the Russian for a shop assistant.
\end{quote}

\subsection*{You'll want to know: \ru{взвешивать}}
\textbf{\ru{взвешивать}} (\emph{vzvéshivat}) — \textquotedblleft{}to weigh\textquotedblright{}.

\begin{quote}
\ru{Продавщица} \ru{взвешивает} \ru{яблоки}.
\end{quote}

\begin{quote}
\emph{Prodavshchítsa vzvéshivayet yábloki.}
\end{quote}

\begin{quote}
The shop assistant weighs the apples.
\end{quote}

Imperfective; perfective \emph{\ru{взвесить}} (\emph{\ru{взвешу}, \ru{взвесишь}}). \emph{\ru{Взвесить} \ru{решение}} — to weigh up a decision.

\begin{cousinweb}
\emph{\ru{взвешивать}} is built from \emph{\ru{вз}-} and \emph{-\ru{вешивать}}, from \emph{\ru{вес}} 'weight'. EtymDB lists the verb without a deeper source.
\end{cousinweb}

\begin{grammarlens}[title={What you've built}]
One more word for telling, explaining and describing at B1.
\end{grammarlens}

\subsection*{Your turn}
\begin{itemize}
\raggedright
  \item \emph{Say it:} \emph{vzvéshivat}
  \item \emph{Say it:} \emph{vzvéshivat}, once more
  \item \emph{Say it:} a sentence of your own with \emph{vzvéshivat}
  \item \emph{Recall:} say the Russian for a counter, then the Russian for a shop assistant, then say \emph{vzvéshivat} again
\end{itemize}

\begin{tcolorbox}[breakable,colback=teal!4,colframe=teal!35!black,title={Before you move on}]
What does \ru{взвешивать} mean? (\textquotedblleft{}To weigh\textquotedblright{}.) Say it once more.
\end{tcolorbox}
